% ID: SEL-D0021-OB-E03
% ESTADO: cotejada. Segunda pasada visual completa y borrador compilado revisados.
% Fuente: Selectividad/matematicas_aplicadas/reserva_a_EXAMEN 3-ANDALUCÍA.pdf
% SHA-256: 0b99df3e2c9a7af0a074136b72c3d257ad9241ccc79fed32894de424481b6be2
% Página: 2; opción B; ejercicio original 3.

Para superar una asignatura un estudiante hace un examen teórico y otro práctico. La probabilidad de que apruebe el examen teórico es $0.8$, la de que apruebe el examen práctico es $0.6$ y la de que apruebe ambos es $0.5$.
\begin{enumerate}[label=\alph*)]
\item \textbf{(1 punto)} \textquestiondown{}Cuál es la probabilidad de que apruebe al menos uno de los dos exámenes?
\item \textbf{(1 punto)} \textquestiondown{}Cuál es la probabilidad de que apruebe el examen práctico en caso de no haber aprobado el examen teórico?
\item \textbf{(0.5 puntos)} \textquestiondown{}Son independientes los sucesos ``aprobar el examen teórico'' y ``aprobar el examen práctico''?
\end{enumerate}
